\clearpage\thispagestyle{empty}
\pdfbookmark[0]{Quellen zu Komet (31. Januar 1814)}{comet1814jan-sources}
\begin{center}\textbf{Quellen zu Komet (31. Januar 1814)}\end{center}
\noindent Diese KI-Übertragung ins heutige Deutsch folgt dem Bericht mit der Überschrift Comet in den \href{https://gdz.sub.uni-goettingen.de/id/PPN236007467}{\textit{Werken}, Band~VI} von Gauß (Göttingen, 1874), S.~373--375. Seine Vorlage ist der Bericht unter Göttingen in den \textit{Göttingischen gelehrten Anzeigen}, Stück~19 (31. Januar 1814), S.~\href{https://api.digitale-sammlungen.de/iiif/image/v2/bsb10538685_00191/full/1600,/0/default.jpg}{185}, \href{https://api.digitale-sammlungen.de/iiif/image/v2/bsb10538685_00192/full/1600,/0/default.jpg}{186}, \href{https://api.digitale-sammlungen.de/iiif/image/v2/bsb10538685_00193/full/1600,/0/default.jpg}{187}, \href{https://api.digitale-sammlungen.de/iiif/image/v2/bsb10538685_00194/full/1600,/0/default.jpg}{188}, \href{https://api.digitale-sammlungen.de/iiif/image/v2/bsb10538685_00195/full/1600,/0/default.jpg}{189}. Die Links führen zum digitalisierten Exemplar der Bayerischen Staatsbibliothek. Der Kometenbericht endet auf S.~188. Im selben Heft folgen ein Bericht über die Präzession und der gesonderte Nachtrag zur Juno.
\par\smallskip\noindent Die Werkausgabe schreibt den Namen des Beobachters Wisniwsky und gibt die aus Bessels zwei Beobachtungen bestimmte Rektaszension des Vergleichssterns IV als 322 44 56.1 an. Im Erstdruck stehen Wisniewsky und 322 44 36,1. Die beiden Doppelpunkte in den Zahlentabellen folgen der Werkausgabe. KI-Übertragung und KI-Satz.
\clearpage
